% !TEX root = ../main.tex
\begin{ouverture}[suites]{2}{Suites numériques}{Partie 1 · Algèbre}
{Termes\enspace·\enspace Récurrence\enspace·\enspace Suites arithmétiques et géométriques\enspace·\enspace Sommes}
\objectif{Calculer les termes d'une suite explicite ou récurrente.}
\objectif{Reconnaître et utiliser une suite arithmétique ou géométrique.}
\objectif{Étudier un sens de variation et calculer une somme.}
\objectif{Modéliser une évolution et programmer un seuil.}
\prerequis{Puissances, pourcentages et calcul littéral.}
\prerequis{Notions de variable, de boucle et de liste en Python.}
\end{ouverture}
\essentiel
\essSection{1}{Définir une suite et lire ses indices}
\begin{definition}
Une suite associe un nombre $u_n$ à chaque entier $n$ d'un ensemble d'indices.
On la note $(u_n)$ ou $(u(n))$ ; le terme de rang $n$ se note $u_n$ ou $u(n)$.
Une formule explicite donne directement $u_n$ ; une relation de récurrence
donne un terme à partir du précédent et nécessite un terme initial.
\end{definition}
\begin{exemple}
Si $u_n=2n+1$ pour $n\geq0$, alors $u_0=1$, $u_1=3$, $u_2=5$.
Si $v_0=1$ et $v_{n+1}=2v_n+1$, alors $v_1=3$, $v_2=7$, $v_3=15$.
\end{exemple}
\begin{attention}
$u_{n+1}$ est le terme de rang $n+1$, tandis que $u_n+1$ signifie « le terme
de rang $n$ auquel on ajoute 1 ». De $u_0$ à $u_n$, il y a $n+1$ termes.
\end{attention}
\essSection{2}{Étudier les variations}
\begin{propriete}
Si $u_{n+1}-u_n\geq0$ pour tout $n$, la suite est croissante ; si cette
différence est négative ou nulle, elle est décroissante. Pour une suite à
termes strictement positifs, on peut aussi comparer $\dfrac{u_{n+1}}{u_n}$ à 1.
\end{propriete}
\begin{center}\begin{tikzpicture}\begin{axis}[otLivreGraphe,width=108mm,height=54mm,xmin=-0.4,xmax=6.5,ymin=-0.5,ymax=14,xtick={0,1,2,3,4,5,6},ytick={2,6,10,14},xlabel={$n$},ylabel={$u_n$}]
\addplot[only marks,otNavy,mark=*,mark size=2pt]coordinates{(0,1)(1,3)(2,5)(3,7)(4,9)(5,11)(6,13)};
\end{axis}\end{tikzpicture}\par\legende{Les points $(n;u_n)$ de $u_n=2n+1$ : les indices sont entiers.}\end{center}
\clearpage
\essSection{3}{Les deux modèles fondamentaux}
\begin{definition}[Suite arithmétique]
$u_{n+1}=u_n+r$, où $r$ est constant. Alors $u_n=u_0+nr$, ou plus généralement
$u_n=u_p+(n-p)r$. Un accroissement constant correspond à ce modèle.
Les points $(n;u_n)$ sont alignés sur la droite de la fonction affine $x\mapsto u_0+rx$.
\end{definition}
\begin{definition}[Suite géométrique]
$v_{n+1}=qv_n$, où $q$ est constant. Alors $v_n=v_0q^n$, ou
$v_n=v_pq^{n-p}$. Un taux constant $t$ correspond à une raison $q=1+t$.
\end{definition}
\begin{exemple}
Ajouter 30 chaque mois donne une suite arithmétique de raison 30.
Augmenter de 3\,\% chaque année donne une suite géométrique de raison $1{,}03$.
Une diminution de 3\,\% correspond à $0{,}97$.
\end{exemple}
\essSection{4}{Calculer une somme sans oublier les bornes}
\begin{theoreme}
Pour une suite arithmétique, la somme de termes consécutifs est
\[\text{nombre de termes}\times\frac{\text{premier terme}+\text{dernier terme}}2.\]
Pour $q\neq1$, $1+q+\cdots+q^n=\dfrac{1-q^{n+1}}{1-q}$.
Pour $q=1$, la somme vaut $n+1$.
\end{theoreme}
\begin{savoirfaire}[Chercher un seuil]
Écrire une boucle qui actualise le terme et son indice tant que le seuil n'est
pas atteint. Vérifier le terme retenu \emph{et son prédécesseur} pour démontrer
qu'il s'agit du premier rang, après avoir établi le sens de variation.
\end{savoirfaire}
\begin{attention}
Des premiers termes suggèrent un comportement ; ils ne démontrent pas une
formule générale. Une suite définie par récurrence n'est pas nécessairement
arithmétique ou géométrique. L'intuition d'une limite ne remplace pas un calcul de seuil.
\end{attention}
\clearpage
\essSection{5}{Observer un comportement à long terme}
\begin{definition}[Une première intuition de limite]
On peut observer qu'une suite se rapproche d'un nombre, augmente au-delà
de tout seuil fixé, ou ne se rapproche d'aucune valeur unique. Cette approche
reste intuitive en Première : aucune définition formelle de limite ni
théorème de convergence n'est nécessaire ici.
\end{definition}
\begin{exemple}
Pour $u_n=1/(n+1)$, les valeurs deviennent proches de 0 tout en restant positives.
Pour $v_n=n^2$, les valeurs dépassent tout seuil fixé si l'on va assez loin.
Pour $w_n=(-1)^n$, les valeurs alternent entre 1 et $-1$ : elles ne se
rapprochent pas d'un seul nombre.
\end{exemple}
\begin{center}\begin{tikzpicture}\begin{axis}[otLivreGraphe,width=110mm,height=61mm,xmin=-.4,xmax=8.5,ymin=-1.4,ymax=1.6,xtick={0,1,2,3,4,5,6,7,8},ytick={-1,0,1},xlabel={$n$},ylabel={valeur}]
\addplot[only marks,otNavy,mark=*]coordinates{(0,1)(1,.5)(2,.333333)(3,.25)(4,.2)(5,.166667)(6,.142857)(7,.125)(8,.111111)};
\addplot[only marks,otGreenDark,mark=square*]coordinates{(0,1)(1,-1)(2,1)(3,-1)(4,1)(5,-1)(6,1)(7,-1)(8,1)};
\end{axis}\end{tikzpicture}\par
\legende{Points ronds : $1/(n+1)$. Carrés : $(-1)^n$.}
\end{center}
\begin{attention}
Une suite croissante n'est pas forcément proche d'un nombre fini : $n^2$
est croissante et grandit sans borne. Un graphique de quelques termes aide
à conjecturer ; la formule de la suite reste indispensable à l'analyse.
\end{attention}
\exercices

\rubrique{Découverte}
\begin{blocLivre}{ex-2-1}
\begin{exercice}[1]{Substituer un indice}
On pose $u_n=n^2-3n+2$ pour $n\geq0$.
\begin{questions}\item Calculer $u_0$, $u_1$, $u_2$ et $u_5$.
\item Exprimer $u_{n+1}$, puis calculer $u_{n+1}-u_n$.
\end{questions}
\end{exercice}
\end{blocLivre}
\begin{blocLivre}{ex-2-2}
\begin{exercice}[1]{Calculer de proche en proche}
$v_0=2$ et $v_{n+1}=3v_n-1$.
\begin{questions}\item Calculer $v_1$, $v_2$, $v_3$ et $v_4$.
\item Les trois premières différences sont-elles constantes ? Et les rapports ?
\end{questions}
\end{exercice}
\end{blocLivre}
\begin{blocLivre}{ex-2-3}
\begin{exercice}[1]{Une suite arithmétique}
$u_0=7$ et la raison est $-3$.
\begin{questions}\item Donner $u_n$, $u_8$ et le sens de variation.
\item Résoudre $u_n\leq-20$ pour $n$ entier naturel.
\end{questions}
\end{exercice}
\end{blocLivre}
\begin{blocLivre}{ex-2-4}
\begin{exercice}[1]{Une suite géométrique}
$v_0=5$ et la raison est 2.
\begin{questions}\item Donner $v_n$, $v_4$ et $v_6$.
\item Calculer $v_0+v_1+v_2+v_3+v_4$ de deux façons.
\end{questions}
\end{exercice}
\end{blocLivre}
\begin{blocLivre}{ex-2-5}
\begin{exercice}[1]{Du pourcentage à la formule}
Un appareil coûte initialement 800 euros. Sa valeur diminue de 15\,\% par an.
\begin{questions}\item Donner la relation entre $V_{n+1}$ et $V_n$ puis $V_n$ en fonction de $n$.
\item Calculer la valeur après deux ans. La baisse globale est-elle de 30\,\% ?
\end{questions}
\end{exercice}
\end{blocLivre}
\clearpage
\rubrique{Entraînement}
\begin{blocLivre}{ex-2-6}
\begin{exercice}[2]{Retrouver les paramètres}
\begin{questions}
\item Une suite arithmétique vérifie $u_3=11$ et $u_8=31$. Déterminer sa raison, $u_0$ et $u_n$.
\item Une suite géométrique de raison positive vérifie $v_2=12$ et $v_5=96$.
Déterminer sa raison, $v_0$ et $v_n$.
\end{questions}
\end{exercice}
\end{blocLivre}
\begin{blocLivre}{ex-2-7}
\begin{exercice}[2]{Compter les termes d'une somme}
\begin{questions}
\item Calculer $4+7+10+\cdots+61$ en déterminant d'abord le nombre de termes.
\item Calculer $3+6+12+\cdots+384$.
\item Calculer $\sum_{k=3}^{10}(2k+1)$.
\end{questions}
\end{exercice}
\end{blocLivre}
\begin{blocLivre}{ex-2-8}
\begin{exercice}[2]{Comparer des variations}
Pour $n\geq0$, on pose $u_n=n^2-4n$ et $v_n=\dfrac{n+1}{n+2}$.
\begin{questions}
\item Calculer $u_{n+1}-u_n$. La suite $(u_n)$ est-elle monotone sur tous les indices ?
\item Montrer que $v_{n+1}-v_n=\dfrac1{(n+2)(n+3)}$ et conclure.
\item Écrire $v_n=1-\dfrac1{n+2}$. Quel comportement semble se dessiner lorsque $n$ grandit ?
\end{questions}
\end{exercice}
\end{blocLivre}
\begin{blocLivre}{ex-2-9}
\begin{exercice}[2]{Lire puis justifier un nuage de points}
Les premiers termes d'une suite sont représentés ci-dessous.
\begin{center}\begin{tikzpicture}\begin{axis}[otLivreGraphe,width=105mm,height=51mm,xmin=-0.4,xmax=5.5,ymin=-0.5,ymax=18,xtick={0,1,2,3,4,5},ytick={2,4,8,16},xlabel={$n$},ylabel={$u_n$}]
\addplot[only marks,otGreenDark,mark=*,mark size=2.3pt]coordinates{(0,16)(1,8)(2,4)(3,2)(4,1)(5,0.5)};
\end{axis}\end{tikzpicture}\end{center}
\begin{questions}\item Conjecturer la nature de la suite et une formule explicite.
\item On précise $u_{n+1}=\frac12u_n$ et $u_0=16$. Justifier la formule et les variations.
\item Déterminer le premier rang pour lequel $u_n<0{,}1$.
\end{questions}
\end{exercice}
\end{blocLivre}
\begin{blocLivre}{ex-2-10}
\begin{exercice}[2]{Comprendre une boucle Python}
\begin{python}
u = 3
termes = [u]
for n in range(4):
    u = 2*u + 1
    termes.append(u)
\end{python}
\begin{questions}\item Donner la liste finale \code{termes}.
\item Quelle relation de récurrence est programmée ? Quel est l'indice du dernier terme ?
\item Adapter le programme pour obtenir les termes de $u_0$ à $u_{10}$.
\end{questions}
\end{exercice}
\end{blocLivre}
\begin{blocLivre}{ex-2-11}
\begin{exercice}[2]{Un seuil d'épargne}
$C_n=1000\times1{,}05^n$ représente un capital après $n$ années.
\begin{questions}\item Justifier que $(C_n)$ est strictement croissante.
\item Écrire une boucle qui détermine la première année où le capital atteint 1500 euros.
\item Exécuter la boucle ou faire des essais, puis vérifier les deux termes qui encadrent le seuil.
\end{questions}
\end{exercice}
\end{blocLivre}
\clearpage
\rubrique{Problèmes}
\begin{blocLivre}{ex-2-12}
\begin{exercice}[3]{Des rangées de sièges}
Une salle a 20 rangées. La première comporte 18 sièges ; chaque rangée suivante
en comporte deux de plus. On note $u_n$ le nombre de sièges de la rangée $n$, pour $1\leq n\leq20$.
\begin{questions}
\item Exprimer $u_n$ et calculer le nombre de sièges de la dernière rangée.
\item Déterminer la capacité totale de la salle.
\item Les dix dernières rangées sont vendues 12 euros la place, les dix premières 9 euros.
Calculer la recette si toutes les places sont vendues.
\item On ajoute des rangées en poursuivant la règle. Quel nombre total minimal de rangées
faut-il pour atteindre au moins 1000 places ?
\end{questions}
\end{exercice}
\end{blocLivre}
\begin{blocLivre}{ex-2-13}
\begin{exercice}[3]{Un traitement répété}
Après chaque heure, l'organisme conserve 80\,\% de la quantité présente, puis
reçoit une nouvelle dose de 6 mg. Juste après la première injection, $u_0=10$ mg.
On mesure $u_n$ juste après la nouvelle dose reçue à l'heure $n$.
\begin{questions}
\item Justifier $u_{n+1}=0{,}8u_n+6$, puis calculer $u_1$ et $u_2$.
\item Poser $v_n=u_n-30$. Montrer que $v_{n+1}=0{,}8v_n$.
\item En déduire $u_n=30-20\times0{,}8^n$ et justifier que $(u_n)$ est croissante et reste inférieure à 30.
\item Avec des essais ou Python, déterminer le premier rang où $u_n\geq28$.
\end{questions}
\label{figlivre-r02-13-cycle-ex}
\figureLivre{r02-13-cycle}

\end{exercice}
\end{blocLivre}
\begin{blocLivre}{ex-2-14}
\begin{exercice}[3]{Deux contrats de rémunération}
La première année, le salaire annuel vaut 24\,000 euros. Le contrat A prévoit
une augmentation de 1000 euros chaque année ; le contrat B, une hausse de 4\,\%.
On note $A_n$ et $B_n$ les salaires de l'année $n+1$ ; ainsi $A_0=B_0=24000$.
\begin{questions}
\item Donner $A_n$ et $B_n$. Comparer les salaires de la sixième année.
\item Calculer le total perçu pendant les dix premières années avec chaque contrat,
puis comparer. Arrondir les montants au centime.
\item Utiliser un tableau ou un programme pour déterminer le premier indice $n\geq1$
où $B_n>A_n$. Vérifier tous les indices antérieurs concernés.
\end{questions}
\end{exercice}
\end{blocLivre}
\begin{blocLivre}{ex-2-15}
\begin{exercice}[3]{Des carrés construits avec des allumettes}
On aligne $n$ carrés unitaires, deux carrés voisins partageant un côté.
\begin{center}\begin{tikzpicture}[x=1.1cm,y=1.1cm]
\foreach \i in {0,1,2,3}{\draw[otNavy,line width=1.1pt](\i,0)rectangle(\i+1,1);}
\node[below,font=\small]at(2,-0.1){Exemple : quatre carrés alignés.};
\end{tikzpicture}\end{center}
\begin{questions}
\item Donner le nombre $a_n$ d'allumettes nécessaires et expliquer la relation de récurrence.
\item Combien de carrés peut-on aligner avec 100 allumettes, sans les casser ?
\item On construit séparément une ligne de 1 carré, puis de 2, jusqu'à $N$ carrés.
Montrer que le nombre total d'allumettes est $T_N=\dfrac{3N^2+5N}{2}$.
\item Déterminer le plus grand $N$ permettant ces constructions avec 1000 allumettes.
\end{questions}
\end{exercice}
\end{blocLivre}
